% !TeX root = ../papers/nc-p3.tex
\section{一辺の二変数商の余次元と残余六次曲線}
\label{sec:binary-incidence-sextic}
\begin{theorem}\label{thm:binary-incidence-codimension}
純関係を含まない$\operatorname{Gr}(6,16)$の開部分で、第一矢を三次元に割る
二変数形式の商の軌跡の閉包は既約で余次元13である。その一般点で核直線は一意である。
関係五平面内の階数二以下の軌跡は、被約支持として次数5の有理正規曲線と
次数15の曲線の和で、後者の正規化は滑らかな平面六次曲線となる。
\end{theorem}
\begin{proof}
核直線の選択は三次元、その商におけるHankel六平面の選択は
$\operatorname{PGL}_3\times\operatorname{PGL}_4$の軌道である。
階数一曲線の二射影が完全線型系で同じ$\mathbb P^1$を埋め込むため、
安定群は対角的な$\operatorname{PGL}_2$であり、軌道次元は$8+15-3=20$である。
分裂を一つ選ぶと持ち上げは$L:k^6\to k^4$で与えられ、24次元である。
従って印付き全空間は既約で次元47となる。一般の持ち上げは純関係を含まない。
実際、Hankel行列$H=(h_{i+j})$の階数一点は一次元しか動かず、
付加する行$Lh$がその行直線に入ることは各点で余次元3である。

階数二以下の四行行列を求める。$B_j(q)$を二次式$q=(q_0,q_1,q_2)$の係数を
一列ずつずらした$(j+1)\times(j+3)$行列とする。
$H$の階数が2なら左核の$[q]\in\mathbb P^2$は一意で$B_3(q)h=0$である。
$H$の二次元の行空間は$\ker B_1(q)$であり、両空間の次元と漸化式から等号となる。
従って付加した行もこの空間に入るための条件は
\[
 M_L(q)h=0,\qquad M_L(q)=\binom{B_3(q)}{B_1(q)L}.
\]
二つの$B_j$は常に全行階数である。$\det M_L$は平面六次式で、
\path{computations/verify_binary_incidence_sextic.py}の整数標本では
素数1009で三つの偏微分に共通射影零点がない。固有性により標数零の一般の
$L$でも滑らかで、従って幾何学的に整である。この場合六次上の$M_L$は常に
階数5なので、核直線が正則に$[h]$を与える。その一般点で$H$は階数2である。
実際、階数一のパラメータ$t$に対し$B_3(q)h(t)=0$は$q(t)=0$となり、
この旗$(t,q)$は二次元である。さらに$B_1(q)Lh(t)=0$は各旗で$L$に
独立な二条件を課すので、一般の$L$では旗の解は高々有限である。
よって$q$の一意性により残余曲線はこの六次と双有理であり、種数は10である。
$H$が階数1ならば付加後も常に階数二以下で、これが五次曲線を与える。

両成分の一般点の重複度は1である。$H$の階数2の開部分では可逆な小行列を用いて
その左核をスキームとして消去でき、上の滑らかな六次の核のグラフとなる。
五次曲線の一般点は二変数の座標で$h=(1,0,\ldots,0)$としてよい。
付加行の末尾三成分を$(u,v,w)$とすると、階数二以下の一次条件は
$(\dot h_2,\dot h_3,\dot h_4)$と$(\dot h_3,\dot h_4,\dot h_5)$がともに
$(u,v,w)$の倍数になることである。一般の$L$では$uw-v^2\ne0$なので、
四つの変分は全て零となり、五次の接方向だけが残る。
四行四列の階数二以下の射影多様体の次数は20である。例えば期待余次元4での
退化軌跡の公式は$ c_2^2-c_1c_3$で、$c_i$は$(1-H)^{-4}$の係数、
従って$100-80=20$である。ここで五平面との交叉は上記の通り期待次元1である。
よって残る次数は15となる。種数10の成分は有理曲線でなく、
前命題の検出する五次曲線は唯一である。その列像から核を復元できるので、
印を忘れても次元47である。Grassmann多様体の次元60との差が主張を与える。
\end{proof}
